\begin{explanationbox}
	\textbf{\textcolor{CExplanation}{一句话直觉}}

	直线的平行与垂直可以通过它们的斜率关系直接判定：平行时斜率相等，垂直时斜率乘积为 $-1$。

	\medskip
	\textbf{\textcolor{CExplanation}{核心拆解}}
	\begin{description}[style=nextline, leftmargin=2.8em, labelsep=0.5em, itemsep=0.45em]
		\item[\textbf{要点一（斜率的意义）：}] 斜率 $k$ 是倾斜角 $\alpha$ 的正切值，即 $k = \tan\alpha$，它反映了直线的倾斜程度。
		\item[\textbf{要点二（平行条件）：}] 若 $l_1 \parallel l_2$，则它们的倾斜角相等（或相差 $\pi$），从而 $\tan\alpha_1 = \tan\alpha_2$，即 $k_1 = k_2$；反之亦然。
		\item[\textbf{要点三（垂直条件）：}] 若 $l_1 \perp l_2$ 且两直线斜率均存在，则倾斜角相差 $\frac{\pi}{2}$，由 $\tan(\alpha+\frac{\pi}{2}) = -\cot\alpha$ 可推出 $k_1k_2 = -1$。
	\end{description}

	\medskip
	\textbf{\textcolor{CExplanation}{几何本质（方向向量关系）}}

	平行 $\Longleftrightarrow$ 方向向量（即 $(1,k_1)$ 与 $(1,k_2)$）成比例；垂直 $\Longleftrightarrow$ 方向向量点积为零，即 $1\cdot1 + k_1k_2 = 0$。

	\medskip
	\textbf{\textcolor{CExplanation}{代数意义（解析化条件）}}

	将位置关系转化为斜率等量关系，是解析几何中“以数论形”的典型代表。只需计算斜率，不需画图即可判断。

	\medskip
	\textbf{\textcolor{CExplanation}{考点价值（常见考法）}}
	\begin{itemize}[leftmargin=*, itemsep=0.5em, topsep=0.4em]
		\item \textbf{考法一（参数求解）：} 根据平行或垂直条件列方程，求直线方程中的待定系数（如 $a$ 的值）。
		\item \textbf{考法二（关系判定）：} 直接比较斜率和斜率乘积，判断两直线位置关系。
	\end{itemize}

	\medskip
	\textbf{\textcolor{CExplanation}{顿悟点}}

	\textbf{垂直的条件是 $k_1k_2=-1$，不是 $k_1=k_2$；平行时还需检查是否重合（截距是否相等）。}

	\medskip
	\textbf{\textcolor{CExplanation}{使用场景}}
	\begin{itemize}[leftmargin=*, itemsep=0.5em, topsep=0.4em]
		\item \textbf{场景一（图形性质求解）：} 求三角形的高线、中线、中垂线方程时，使用垂直条件求斜率。
		\item \textbf{场景二（参数范围讨论）：} 已知两直线平行或垂直，求参数取值范围，或讨论位置关系。
	\end{itemize}
\end{explanationbox}
